\documentclass[../../main.tex]{subfiles}
\begin{document}

\chapter{Evaluation methodologies and benchmarks for code LLMs}
\label{ch:evaluation_metrics_benchmarks_se}

\section{Chapter overview and learning outcomes}
Evaluating code-generating models requires fundamentally different paradigms than evaluating natural language models. A code snippet with a single misplaced semicolon or flipped boolean condition achieves high surface text similarity while being completely broken. This chapter explores syntactic metrics, derives execution-based metrics ($\PassAtK{k}$), and analyzes standard software engineering benchmarks.

Upon completing this chapter, you will be able to:
\begin{itemize}
    \item Explain the failure modes of surface NLP metrics (BLEU, ROUGE) when applied to source code.
    \item Formulate the multi-aspect components of the CodeBLEU metric.
    \item Derive the unbiased combinatorial estimator of the $\PassAtK{k}$ metric \citep{chen2021codex}.
    \item Contrast synthetic standalone benchmarks (HumanEval, MBPP) with repository-level benchmarks (SWE-bench \citep{jimenez2024swebench}).
    \item Identify benchmark contamination and design decontamination filtering pipelines.
\end{itemize}

\section{Surface and syntactic metrics: CodeBLEU}
Standard BLEU measures $n$-gram lexical overlap between generated code and reference implementations. However, renaming variable \texttt{counter} to \texttt{idx} severely penalizes BLEU despite having zero semantic impact.

\begin{definition}{CodeBLEU metric}{codebleu_metric}
CodeBLEU combines four distinct structural and semantic dimensions via a weighted linear combination:
\begin{equation}
\label{eq:codebleu}
\CodeBLEU = \alpha \cdot \text{BLEU} + \beta \cdot \text{BLEU}_{\text{weight}} + \gamma \cdot \text{Match}_{\text{AST}} + \delta \cdot \text{Match}_{\text{DF}},
\end{equation}
where $\text{Match}_{\text{AST}}$ evaluates tree edit similarity on parsed Abstract Syntax Trees, $\text{Match}_{\text{DF}}$ evaluates overlap on Data Flow def-use graphs, and weights typically sum to $1$ ($\alpha = \beta = \gamma = \delta = 0.25$).
\end{definition}

\section{Execution-based evaluation: the Pass@k metric}
Because code is executable, the ultimate measure of functional correctness is running the generated code against a unit test suite.

\begin{theorem}{Unbiased combinatorial estimator of Pass@k}{pass_at_k}
Suppose that for each programming problem, we sample $n$ candidate code solutions from the model ($n \ge k$), and evaluate each against unit tests. Let $c$ be the number of generated samples that pass all test cases ($0 \le c \le n$).
The naive estimator $1 - (1 - c/n)^k$ suffers from severe bias for small $n$. Chen et al. \citep{chen2021codex} derived the uniformly minimum-variance unbiased estimator of $\PassAtK{k}$:
\begin{equation}
\label{eq:pass_at_k_unbiased}
\PassAtK{k} = \E_{\text{problems}} \left[ 1 - \frac{\binom{n - c}{k}}{\binom{n}{k}} \right] = \E_{\text{problems}} \left[ 1 - \frac{\prod_{i=0}^{k-1} (n - c - i)}{\prod_{i=0}^{k-1} (n - i)} \right],
\end{equation}
with the convention that if $n - c < k$, the ratio $\frac{\binom{n-c}{k}}{\binom{n}{k}} = 0$, hence $\PassAtK{k} = 1.0$.
\end{theorem}

\begin{examinsight}[Exam derivation: combinatorial probability of Pass@k]
Oral and written exams frequently ask students to derive Equation~\eqref{eq:pass_at_k_unbiased}:
\begin{enumerate}
    \item The total number of ways to choose a subset of $k$ samples from $n$ generated candidates is $\binom{n}{k}$.
    \item A problem \textbf{fails} under $\PassAtK{k}$ if and only if \emph{all} $k$ selected samples are incorrect.
    \item The number of incorrect samples is $n - c$. The number of ways to choose $k$ incorrect samples is $\binom{n-c}{k}$.
    \item The probability that all $k$ samples fail is therefore $\frac{\binom{n-c}{k}}{\binom{n}{k}}$.
    \item The probability that \textbf{at least one} sample passes is $1 - \frac{\binom{n-c}{k}}{\binom{n}{k}}$.
\end{enumerate}
This derivation requires zero approximations and produces completely unbiased estimates.
\end{examinsight}

\begin{deepdive}[Benchmark contamination and decontamination]
Because modern LLMs pre-train on terabytes of public GitHub repositories, test benchmarks such as HumanEval or LeetCode problems may have been present in pre-training corpora. Evaluating on contaminated data yields misleadingly inflated $\PassAtK{k}$ scores. State-of-the-art evaluation requires $n$-gram decontamination filtering, syntactic perturbation, or dynamic test execution on freshly committed repositories.
\end{deepdive}

\end{document}
